% !TEX root = ../report.tex

\section{Discussion}
\label{sec:discussion}

The three pieces of evidence answer the paper's question---does diversification fail exactly when it is needed?---in a way that is more precise than a simple ``yes''. Read together, they separate the diversification that is genuinely lost in stress from the diversification that only \emph{appears} to be lost because volatility has risen, and they locate where each kind sits.

\paragraph{Diversification does weaken in realized, experienced terms.} The raw regime-conditional correlations and the principal-component concentration both move in the fragile direction. In the international-equity panel the whole dependence structure shifts up---the average pairwise correlation rises from $0.74$ to $0.87$ and the leading component absorbs $89\%$ of variance in stress, leaving barely more than one effective bet across six markets. In the cross-asset panel the average correlation does not rise, yet the portfolio still loses about one-fifth of an independent bet as the risk-on block converges onto a single direction. Every universe is more concentrated in stress. This is the comovement an investor actually experiences, and on this measure diversification is regime-dependent: a covariance matrix estimated over a full sample understates the joint risk that materializes in a crisis, exactly the instability of the optimization inputs anticipated in the introduction.

\paragraph{The increase is volatility-driven interdependence, and it is concentrated in equities.} The Forbes--Rigobon adjustment locates the breakdown rather than dissolving it. Once the stress correlations are deflated for the higher variance of the common factor, \emph{no} pair in any universe survives as a robust breakdown: in the language of \citet{forbesrigobon2002}, the evidence points to heightened \emph{interdependence} in a more volatile state rather than \emph{contagion}, a structural break in the dependence relationship. This relabelling does not rescue the investor's diversification---the experienced comovement is unchanged---but the conservative pair-level correction exposes a sharp asymmetry that the average numbers hide. In the cross-asset panel nineteen of the twenty-eight non-SPY pairs are genuinely resilient, their correlation \emph{falling} in stress, led by Treasuries and gold; the international-equity panel contains no resilient pair and supplies almost all of the volatility-bias-sensitive increases. The honest headline is therefore qualified and specific: diversification weakens in stress, but largely because equity-block variance rises rather than because dependence structurally changes, and the cross-asset hedges keep---indeed strengthen---their diversifying role exactly when it is needed.

\paragraph{A portfolio-management reading.} These results are diagnostics for portfolio construction, not allocation rules. The Panel~B average shows that adding non-U.S.\ equity markets does not create independent crisis exposures: the international-equity block tightens almost uniformly, so geographic breadth buys little diversification in precisely the states where it is needed. The Panel~A average shows that cross-asset diversification is more robust---but not because every non-equity sleeve turns defensive. Listed real estate, commodities, oil, and high-yield credit remain part of the risk-on complex, while sovereign duration and gold supply the genuinely distinct stress exposures; the VNQ-linked increases are the clearest case, economically consistent with listed real estate behaving as an equity-like risk asset in stress rather than as a separate real-asset hedge. The implication is qualitative rather than a prescription: stress-robust diversification requires genuinely different risk \emph{drivers}, not merely more assets or more countries, and a candidate diversifier is best judged by whether its correlation holds or falls in the stress regime rather than by its unconditional behaviour.

\paragraph{Limitations.} Several caveats bound these conclusions. First, the stress sample is small: the regime contains only a handful of distinct episodes, so the regime-conditional estimates rest on limited independent variation. Second, the analysis is associational---the Markov regime is an ex-post label, not an exogenous instrument---so none of the results identify a causal contagion mechanism. Third, the ETF histories constrain the balanced window to begin in 2007, and the universe, while economically representative, is deliberately parsimonious. Fourth, the regime itself is model- and rule-dependent, although the agreement between the Markov and VIX-plus-drawdown labels mitigates this. Fifth, the stress regime pools heterogeneous macroeconomic episodes; the descriptive stress-episode split of Appendix Table~\ref{tab:stress_episode_hedges} shows that hedge signs---most notably the stock--Treasury relation---can differ across stress types. Finally, the Forbes--Rigobon correction is exact only under its single-factor, variance-only assumption and depends on the chosen shock variance; the \texttt{market\_spy} and \texttt{pair\_max} variants bracket this choice but do not eliminate it. These limitations qualify the magnitude of the findings, not their direction.
